\section{Introduction}
\label{sec:intro}

Deep model merging has recently emerged as an essential paradigm for multi-task learning and model reuse, offering a training-free pathway to combine multiple task-specific expert neural networks into a single multi-task model \cite{ilharco2022editing, wortsman2022model, yadav2023resolving}. Unlike standard multi-task fine-tuning, which requires joint access to all training datasets and is highly compute-intensive, model merging directly interpolates the weights or task vectors of expert models fine-tuned from a shared pre-trained base model. Simple strategies like uniform weight averaging (Task Arithmetic) \cite{ilharco2022editing} or sign-consensus filtering (TIES-Merging) \cite{yadav2023resolving} provide robust initial multi-task capabilities, but suffer from significant interference when merging task experts with conflicting parameter gradients or disparate scaling behaviors.

To address inter-task interference, adaptive merging methods have been proposed. Notable among these is AdaMerging \cite{yang2024adamerging}, which introduces learnable coefficient parameters for different layers and tasks. Crucially, because model merging assumes no access to the original training datasets, AdaMerging operates at test-time by running unsupervised test-time adaptation (TTA) via entropy minimization on small batches of unlabeled target data. By optimizing layer-wise merging coefficients to minimize prediction entropy, AdaMerging dynamically tailors the joint parameter configuration to the local test-time stream, achieving notable accuracy gains on multi-task benchmarks.

Despite these promising gains, this paper deconstructs the adaptive test-time merging paradigm and exposes two severe, under-studied failure modes that challenge its validity and robustness. 

\subsection{The Overfitting-Optimizer Paradox}
The first critical vulnerability is what we term the \textbf{Overfitting-Optimizer Paradox}. Standard layer-wise optimization (AdaMerging) employs first-order gradient descent (e.g., Adam) or evolutionary strategies (e.g., $1+1$ ES) to adjust coefficients independently across dozens of model layers. Typically, these coefficients are optimized on a highly restricted transductive window (e.g., a single calibration batch of $16$ or $32$ samples per task). 

Our diagnostic audits reveal that optimization of fine-grained layer-wise coefficients overfits heavily to the transductive calibration statistics rather than learning a generalized cross-basin parameter configuration. To empirically expose this, we introduce a \emph{spatial shuffling diagnostic}: after optimizing layer-wise coefficients using standard AdaMerging, we randomly shuffle the optimized coefficients across the model's layers. Remarkably, we observe that the spatially shuffled coefficients retain almost all of the performance gains of the original "localized" coefficients (as shown in our experiments in Table~\ref{tab:main_results}). This indicates that the localized layer-by-layer optimization does not capture genuine spatial feature interactions; rather, it behaves as an unconstrained parameter-drift mechanism that fits the noise of the small calibration batch, suffering from severe transductive overfitting.

\subsection{Sacrificial Task Bias}
The second severe limitation is the \textbf{Sacrificial Task Bias}. In multi-task model merging, task experts are evaluated on domains of vastly different complexities (e.g., comparing highly regularized digit recognition on MNIST to complex, real-world street numbers on SVHN). Because these tasks exhibit vastly different baseline entropy ranges and class capacities, optimizing a single joint entropy loss:
\begin{equation}
\mathcal{L}_{\text{joint}} = \sum_{k=1}^K \mathcal{H}_k(\Lambda)
\end{equation}
results in an extremely biased gradient landscape. Under standard optimization, tasks with low baseline entropy and high complexity are heavily penalized or ignored, as the joint optimizer can minimize the total loss more easily by over-optimizing simpler tasks. Consequently, complex tasks like SVHN are actively "sacrificed," suffering severe degradation in test accuracy compared to naive, uniform Task Arithmetic (as demonstrated in our baseline evaluations where SVHN drops in accuracy under unconstrained evolutionary optimization).

\subsection{Our Proposed Framework: RegCalMerge}
To address both transductive overfitting and sacrificial task bias, we propose \textbf{RegCalMerge} (Calibrated \& Regularized Test-Time Model Merging), a robust, unified framework that can be deployed across a spectrum of adaptation priorities:
\begin{enumerate}
    \item \textbf{Calibrated AdaMerging (CalMerge)} ($\beta=0.0, \gamma=0.0$): A high-performance calibration engine combining \emph{Class-Capacity Normalization} (CCN) and \emph{Scale-Normalized Entropy Weighting} (SNEW). CCN normalizes raw prediction entropy by its maximum theoretical capacity ($\log C_k$, where $C_k$ is the class count) to map disparate entropy values onto a uniform $[0, 1]$ interval. SNEW weights each normalized task entropy by the inverse of its baseline uniform task arithmetic entropy at step $0$. Together, they ensure that complex domains contribute equally to joint gradients, completely resolving the Sacrificial Task Bias. This unregularized calibration engine achieves a state-of-the-art Joint Mean accuracy of \textbf{61.82\%} (outperforming all existing baselines) and raises SVHN accuracy to \textbf{32.03\%} (the highest in the literature).
    \item \textbf{Regularized Calibrated Merging (RegCalMerge)} ($\beta > 0, \gamma > 0$): For settings where parameter safety is paramount, we introduce \emph{Elastic Spatial Regularization} (ESR) as an optional structural stabilizer. ESR applies a \emph{Proximity Penalty} ($\beta$) keeping coefficients near their uniform Task Arithmetic initialization, and a \emph{Spatial Deviation Penalty} ($\gamma$) penalizing the variance of layer-wise coefficients around their task-wise spatial average. This acts as a smooth, adjustable spatial constraint that restricts high-frequency parameter drift and transductive overfitting, trading off a small fraction of local peak accuracy to guarantee global parameter stability.
\end{enumerate}

To rigorously evaluate RegCalMerge, we adopt the philosophy of \textbf{The Empiricist}, running extensive grid sweeps crossing the ESR parameters ($\beta, \gamma$) over multiple random seeds, and evaluating multiple baseline configurations. Our dense hyperparameter sweeps map out a smooth, highly predictable generalization landscape, proving that RegCalMerge acts as a robust, calibration-aware dial that resolves transductive overfitting while maintaining localized representational benefits. Our contributions are summarized as follows:
\begin{itemize}
    \item We expose the \textbf{Overfitting-Optimizer Paradox} in test-time model merging via a novel spatial shuffling diagnostic, proving that unregularized layer-wise adaptive merging overfits to local calibration samples rather than learning stable, localized layer interactions.
    \item We identify and analyze the \textbf{Sacrificial Task Bias}, showing how unnormalized multi-task entropy objectives systematically degrade performance on complex, high-entropy domains like SVHN.
    \item We introduce \textbf{RegCalMerge}, a robust calibration-aware merging framework comprising a core calibration engine (\textbf{CalMerge}) that resolves sacrificial task bias to achieve state-of-the-art performance, and an optional regularization module (\textbf{ESR}) that stabilizes weight-space dynamics.
    \item We present a large-scale, multi-seed empirical sweep over 7 baseline configurations, showing that our calibrated layer-wise formulation out-performs a calibrated spatially-averaged baseline (\textbf{61.82\%} vs \textbf{61.13\%}), thereby confirming that fine-grained layer-wise parameter flexibility is essential but requires calibration to prevent destructive domain bias.
\end{itemize}